\par\addvspace{12pt}\Needspace{9\baselineskip}

\section{数据的获取}

% ITEM: probability-example-027
\solutionitem{probability-example-027}{例6.27}
某校准备调查全校学生的睡眠时间。比较以下方案：在图书馆晚间离馆者中调查；在全校学生名单中简单随机抽取；按年级分层随机抽取。说明各方案的目标总体、可能遗漏的对象，以及分层比例不同时怎样汇总平均值。
\solutionanswer
三种方案原本的目标总体都应是规定时段内的全校学生。晚间图书馆离馆者只能直接代表这部分人，未到图书馆、较早离馆等学生没有被覆盖，睡眠习惯也可能与入馆行为有关。全校完整名单上的简单随机抽样使每个规定大小的子集等可能，还需按统一口径记录睡眠时间并处理未回应。按年级分层可以保证各年级进入样本；若抽样比例不同，应按总体人数比例\(N_h/N\)汇总\(\sum_h(N_h/N)\bar x_h\)，而不是按各层实际抽到的人数比例直接合并。


% ITEM: probability-exercise-18-01
\solutionitem{probability-exercise-18-01}{习题6.18（1）}
某总体分三层，人数为500、300、200，计划抽取100人。给出比例分配方案。若改为各层抽30人，三个样本均值为12、15、18，求总体均值的分层估计。
\solutionanswer
比例分配为\(50,30,20\)人。等额各抽30人后，仍须按总体权重\(0.5,0.3,0.2\)汇总，估计均值\(0.5\times12+0.3\times15+0.2\times18=14.1\)，不是三个样本均值的简单平均\(15\)。


% ITEM: probability-exercise-18-02
\solutionitem{probability-exercise-18-02}{习题6.18（2）}
为了研究某训练方法，从报名者中随机分为训练组与对照组。说明随机分组解决了什么问题；若希望结论适用于所有高中生，还需要考虑哪些取样信息？
\solutionanswer
随机分组使报名者接受哪一种处理不由原有能力、积极性等特征系统决定，为组间效果比较提供依据。要推广到所有高中生，还需考察报名者如何进入研究、覆盖了哪些地区年级、未报名者与报名者可能有什么差异，以及测量和执行是否可比。随机分配处理并未同时完成从全部高中生随机抽样。


% ITEM: probability-exercise-18-03
\solutionitem{probability-exercise-18-03}{习题6.18（3）}
从编号1至20的对象中随机选一个起点\(1,\ldots,5\)，再每隔5个取一个。求每个对象的入选概率，并判断这是否是样本量4的简单随机抽样。
\solutionanswer
可能样本为\(\{r,r+5,r+10,r+15\}\)，\(r=1,\ldots,5\)，每个概率\(1/5\)。每个对象恰属于其中一个样本，入选概率均为\(1/5=4/20\)。但只有五个四元子集可能被抽到，其余子集概率零，故不是样本量4的简单随机抽样。


% ITEM: probability-exercise-18-04
\solutionitem{probability-exercise-18-04}{习题6.18（4）}
（据2022年新高考II卷第19题改编）随机调查某地区100名某疾病患者，将年龄分入\([0,10),[10,20),\ldots,[80,90]\)，得到下表：
\[
\begin{array}{c|rrrrrrrrr}
 \text{组中值（岁）}&5&15&25&35&45&55&65&75&85\\\hline
 \text{频数}&1&2&12&17&23&20&17&6&2
\end{array}
\]
各组以组中值代表，估计患者的平均年龄及年龄在\([20,70)\)内的概率。若全地区患病率为\(0.1\%\)，年龄在\([40,50)\)内的人口占\(16\%\)，估计从这一年龄段随机选一人时其患病的概率，精确到\(0.0001\)。说明最后一问所用的总体与患者样本之间的关系。
\solutionanswer
以组中值加权，估计平均年龄\(\sum_hm_hn_h/100=47.9\)岁。\([20,70)\)的患者频数\(12+17+23+20+17=89\)，概率估计\(0.89\)。令\(D\)为患病，\(A\)为40至50岁，则患者样本给\(P(A\mid D)\approx0.23\)，地区资料给\(P(D)=0.001,P(A)=0.16\)。所以\(P(D\mid A)\approx0.23\times0.001/0.16=0.0014375\)，精确到\(0.0001\)为\(0.0014\)。最后一问的分母来自全地区该年龄段人口，不能将患者样本中该年龄组的频率直接当作患病率。
\par
